\documentclass[10pt]{article}

%%% Doc layout
\usepackage{fullpage} 
\usepackage{booktabs}       % professional-quality tables
\usepackage{graphicx}       % figure graphics
\usepackage{microtype}      % microtypography
\usepackage{parskip}
\usepackage{times}


%% Hyperlinks always black, no weird boxes
\usepackage[hyphens]{url}
\usepackage[colorlinks=true,allcolors=black,pdfborder={0 0 0}]{hyperref}

%%% Math typesetting
\usepackage{amsmath,amssymb}

%%% Write out problem statements in purple, solutions in black
\usepackage{xcolor}
\newcommand{\officialdirections}[1]{{\color{purple} #1}}

%%% Avoid automatic section numbers (we'll provide our own)
\setcounter{secnumdepth}{0}


%% --------------
%% Header
%% --------------
\usepackage{fancyhdr}
\fancyhf{}
\setlength{\headheight}{15pt}
\fancyhead[C]{\ifnum\value{page}=1 Tufts CS 136 - 2026f - CA3 Submission \else \fi}
\fancyfoot[C]{\thepage} % page number
\renewcommand\headrulewidth{0pt}
\pagestyle{fancy}

%% --------------
%% Begin Document
%% --------------
\begin{document}

~\\ %% add vertical space
{\Large{\bf Student Name: TODO}}

~\\ %% add vertical space
{\bf Collaboration Statement:}

Total hours spent: TODO

I consulted the following human, textbook, or AI resources:
\begin{itemize}
\item TODO
\item TODO
\item $\ldots$	
\end{itemize}

By turning this document in, I attest that I have followed the 
\href{https://www.cs.tufts.edu/cs/136/2026f/index.html#collaboration}{[CS 136 Collaboration Policy]}. All solutions represent my own work. No solution text in this document was directly provided by other humans or artificial agents. 

\setcounter{tocdepth}{2}
\tableofcontents

\newpage


\subsection{1a}
\renewcommand{\thefigure}{1a}
 \begin{figure}[!h]
     \centering
     \includegraphics[width=0.32\textwidth]{example-image-a.pdf}
     \includegraphics[width=0.32\textwidth]{example-image-b.pdf}
     \includegraphics[width=0.32\textwidth]{example-image-c.pdf}
     \label{fig:1a}
     \caption{
TODO YOUR CAPTION HERE (See assignment html for instructions)
}%endcaption
 \end{figure}

\vspace{2cm}

\subsection{1b}
\renewcommand{\thefigure}{1b}
 \begin{figure}[!h]
     \centering
     \includegraphics[width=0.32\textwidth]{example-image-a.pdf}
     \includegraphics[width=0.32\textwidth]{example-image-b.pdf}
     \includegraphics[width=0.32\textwidth]{example-image-c.pdf}
     \label{fig:1b}
     \caption{
TODO YOUR CAPTION HERE (See assignment html for instructions)
}%endcaption
 \end{figure}



\subsection{1c}
\officialdirections{If you want good reconstructions, which do you care more about: making sure there are no crowded regions in the latent space, or making sure there are no empty regions? \\
If you want good generations/samples from $p(x)$, which do you care more about: making sure there are no crowded regions in the latent space, or making sure there are no empty regions?
}

TODO YOUR SOLUTION HERE

\newpage 

\subsection{2a}
\officialdirections{Look at the final equation you implemented for the KL divergence term of the loss. Then answer: if we want the KL divergence to be small, what value is it pushing $\sigma$ towards? Does your answer depend on what the encoder network $\mu_\phi(x_n)$ is doing?
}

TODO YOUR SOLUTION HERE

\subsection{2b}
\officialdirections{Unlike 2a, you can answer this question using intuition instead of math. Discuss: if you want the reconstruction loss $\mathbb{E}_{q_\phi(z_n|x_n)}\left[-\log p_\theta(x_n|z_n)\right]$ to be minimized, then do you want $\sigma$ to be small or large? Does your answer depend on what the encoder network $\mu_\phi(x_n)$ is doing?

Combining this with the answer to 2a, explain why $\sigma$ grows at first then shrinks.
}

TODO YOUR SOLUTION HERE
 
\newpage 

\subsection{3a}
\renewcommand{\thefigure}{3a}
 \begin{figure}[!h]
     \centering
     \includegraphics[width=0.32\textwidth]{example-image-a.pdf}
     \includegraphics[width=0.32\textwidth]{example-image-b.pdf}
     \includegraphics[width=0.32\textwidth]{example-image-c.pdf}
     \label{fig:3a}
     \caption{
TODO YOUR CAPTION HERE (See assignment html for instructions)
}%endcaption
 \end{figure}

\vspace{2cm}

\subsection{3b}
\renewcommand{\thefigure}{3b}
 \begin{figure}[!h]
     \centering
     \includegraphics[width=0.32\textwidth]{example-image-a.pdf}
     \includegraphics[width=0.32\textwidth]{example-image-b.pdf}
     \includegraphics[width=0.32\textwidth]{example-image-c.pdf}
     \label{fig:3b}
     \caption{
TODO YOUR CAPTION HERE (See assignment html for instructions)
}%endcaption
 \end{figure}

\newpage 

\subsection{4}
\officialdirections{Using the equation above, compute a lower bound on the amount of probability mass allocated to the 60,000 images in our training dataset, and a lower bound on how much probability mass is being allocated to the 10,000 images in our test dataset (note that what is reported to you in the csv is $\mathcal{L}/N$ because it's per-image). Then discuss:

\begin{itemize}
\item Up to how much of the probability mass is leftover?
\item Does the ELBO (or the log evidence) care about how the remaining probability mass is distributed?
\item Does the quality of generations/samples from our model care about how the remaining probability mass is distributed?
\end{itemize}
}

TODO YOUR SOLUTION HERE

\end{document}
